%% 50-discussion.tex — Discussion

\section{Discussion}
\label{sec:discussion}

\subsection{Known Alignment Gaps}
\label{subsec:gaps}

Several field classes remain outside the current pipeline scope.
These are documented explicitly rather than silently dropped, so that
downstream users can assess whether a given query will be affected:

\begin{itemize}
  \item \textbf{EDM-structural properties.}
        \texttt{edm:isShownAt}, \texttt{edm:aggregatedCHO},
        \texttt{edm:dataProvider}, and related provenance and access properties
        have no DC$\to$RDA route and require a separate direct mapping to
        mocho or PROV-O.
        They constitute 115{,}432 occurrences for \texttt{edm:isShownAt} and
        \texttt{edm:aggregatedCHO} alone (one per record).

  \item \textbf{EDM agent authority properties.}
        \texttt{edm:dateOfBirth}, \texttt{edm:placeOfBirth},
        \texttt{edm:biographicalInformation}, and \texttt{edm:gender} on Agent
        entities require direct mapping to RDA Person/Corporate Body
        properties, not a DC-mediated route.
        This affects all Agent entities in the corpus.

  \item \textbf{SKOS labels.}
        \texttt{skos:prefLabel} and \texttt{skos:altLabel} on Agent, Concept,
        and Place entities have no RDA sub-property in RDA~5.4.9 and are
        excluded.
        A SKOS-to-RDA access-point alignment step is required.

  \item \textbf{Identifier and spatial properties.}
        \texttt{dc:identifier} and \texttt{dcterms:spatial} have no
        RDA sub-property in the DC\,$\to$\,RDA mapping file; both are
        documented as known gaps.

  \item \textbf{A DDB API typo.}
        \texttt{Event.occuredAt} (sic) appears in 102{,}859 records but
        resolves to no IRI (expected \texttt{edm:occurredAt}).
        It is excluded pending confirmation from DDB.
        \todo{File an issue with DDB or cite if already documented.}
\end{itemize}

\subsection{Corpus Scope and Generalizability}

The pilot corpus is drawn from DDB search results for \textit{Goethe} and
\textit{Faust}: it is a thematically focused convenience sample, not a
probability sample of the full DDB corpus.
The corpus covers all seven sectors and all five mediatypes, so the dispatch
tables are populated across all structural categories, but sector proportions
will differ from the full 65M-record corpus.
Claims about fallback rates and W-slot hit rates should be understood as
estimates conditional on this corpus composition.

The three-layer dispatch architecture is designed to generalize to any
DDB-EDM corpus: the dispatch tables are keyed on sector and mediatype codes
that are stable across the DDB.
Extension to the full 65M-record corpus requires re-running the corpus
profiling step and verifying that no new (entity type, JSON key) pairs
emerge that are not covered by the current alignment table.
\todo{Check: are there known sector/mediatype combinations absent from the
  goethe-faust pilot? E.g., sparte003 (Denkmal) may be underrepresented.
  Report sector distribution table if available from corpus analysis.}

\subsection{Reproducibility}

All dispatch tables are versioned alongside the pipeline script.
The alignment table is reproducible: re-running the profiling and alignment
scripts on the same JSONL input produces byte-identical output.
The mocho OWL import chain pins specific ontology versions; the
\texttt{in\_mocho} flag in the alignment table therefore reflects a
specific mocho release \todo{confirm version tag for paper}.
Downstream users consuming the N-Triples output should be aware that a
mocho release update may alter which properties are flagged as \texttt{in\_mocho},
and should re-run the pipeline against the updated import if needed.

\subsection{Relation to Automated Alignment Tools}

The methodology does not replace automated alignment tools; it establishes
the preconditions under which they can be applied reliably.
Once the corpus profiling step reveals that \texttt{dc:type} carries three
distinct semantic regimes, a system such as BERTMap~\cite{he2022bertmap} can
be applied \emph{per stratum} rather than globally, improving the signal it
receives.
The dispatch tables produced in Layer~3 could also serve as training data or
evaluation benchmarks for automated tools on this domain.
